% =====================================================================
%  RAG project: audit, head-to-head benchmark, and resume guidance
%  Version 2 -- supersedes rag_project_audit.tex
%  Compile:  pdflatex rag_project_audit_v2.tex   (run twice for the ToC)
% =====================================================================
\documentclass[11pt,a4paper]{article}

\usepackage[margin=2.1cm]{geometry}
\usepackage{booktabs}
\usepackage{longtable}
\usepackage{array}
\usepackage{xcolor}
\usepackage{enumitem}
\usepackage{listings}
\usepackage[hidelinks]{hyperref}
\usepackage{titlesec}
\usepackage{ragged2e}

\definecolor{supported}{RGB}{0,120,60}
\definecolor{partial}{RGB}{190,120,0}
\definecolor{unsupported}{RGB}{180,30,30}
\definecolor{codebg}{RGB}{245,245,245}
\definecolor{boxbg}{RGB}{248,248,242}

\newcommand{\SUP}{\textcolor{supported}{\textbf{SUPPORTED}}}
\newcommand{\PART}{\textcolor{partial}{\textbf{PARTIAL}}}
\newcommand{\NOT}{\textcolor{unsupported}{\textbf{NOT SUPPORTED}}}

% narrow, ragged-right column type for evidence cells so long file:line
% strings wrap instead of running into the margin
\newcolumntype{R}[1]{>{\RaggedRight\arraybackslash}p{#1}}

\lstdefinestyle{code}{
  basicstyle=\ttfamily\footnotesize,
  backgroundcolor=\color{codebg},
  breaklines=true,
  frame=single,
  framesep=4pt,
  rulecolor=\color{gray!40},
  columns=fullflexible,
  keepspaces=true,
  showstringspaces=false,
  xleftmargin=0pt,
}
\lstset{style=code}

\titleformat{\section}{\large\bfseries}{\thesection}{0.7em}{}
\setlist{itemsep=2pt,parsep=0pt,topsep=3pt}

\title{\textbf{RAG Project: Audit, Head-to-Head Benchmark,\\ and Resume Guidance}\\[3pt]
\large From-Scratch NumPy Pipeline vs.\ LangChain/LangGraph Transcript Assistant}
\author{Parikshit M. \quad (M.Tech Research, CDS, IISc Bengaluru)}
\date{24 September 2026 \quad$\cdot$\quad Version 2}

\begin{document}
\maketitle

% =====================================================================
\section*{Summary}
\addcontentsline{toc}{section}{Summary}
% =====================================================================

\noindent
\fcolorbox{gray!40}{boxbg}{%
\begin{minipage}{0.97\textwidth}
\vspace{4pt}
\textbf{In one paragraph.} This project is two implementations of the same idea,
and the interesting result is that they were finally measured against each other.
The from-scratch pipeline (\texttt{rag\_annotated.py}: NumPy, MiniLM-384d,
800/150-character chunks, dot-product cosine search) and the framework pipeline
(LangChain + LangGraph + Chroma + Groq \texttt{gpt-oss-120b}, 1000/200-character
chunks) were run over the same 52{,}826-character transcript against the same 35
span-labelled questions. At their shipped defaults the from-scratch version
retrieves slightly \emph{better} --- hit@4 of 0.857 versus 0.829 --- and indexes
faster. A chunk-size ablation then shows that this gap is not about NumPy versus
Chroma at all: it is chunk size. Holding chunking fixed at 800/150, the framework
pipeline rises to hit@4 0.886 with a perfect hit@10 of 1.000, its best
configuration measured, and better than the 1000/200 setting the application
currently ships. Retrieval costs 8--10\,ms in both systems, so the vector store is
not the latency story --- the LLM call is. The honest headline for a resume is
therefore not ``framework good'' or ``from-scratch good'' but: \emph{two
implementations, one gold set, and a measured finding that chunk size dominates
both}. Two words from the original resume text remain unusable: the application
is \textbf{served locally, not deployed}, and its two-node graph makes no
decisions, so it is \textbf{not an agent}.
\vspace{4pt}
\end{minipage}}

\tableofcontents
\vspace{1em}

% =====================================================================
\section{What changed since version 1}
% =====================================================================

Version 1 of this audit could not verify the from-scratch claim: the code was not
in the repository, so bullet 1 was marked \emph{not verifiable} and bullet 3 --- the
comparison --- had no numbers behind it at all.

\texttt{rag\_annotated.py} is now in the repository, at
\texttt{DL/Lang\_Chain/Lang\_Chain/rag\_annotated.py}. That closes both gaps:

\begin{itemize}
  \item Bullet 1 is now auditable, and the wording needs one correction
        (Section~\ref{sec:bullet1}): the pipeline is written around a
        \emph{pretrained} sentence-transformers encoder, not a from-scratch
        embedding model.
  \item Bullet 3 is now a measurement rather than a claim
        (Section~\ref{sec:bench}), including a chunk-size ablation that
        identifies \emph{why} one pipeline beats the other.
\end{itemize}

% =====================================================================
\section{The two implementations, side by side}
% =====================================================================

\begin{longtable}{@{}R{3.1cm}R{5.9cm}R{6.1cm}@{}}
\toprule
& \textbf{From-scratch (\texttt{rag\_annotated.py})} &
  \textbf{Framework (\texttt{Transcript\_RAG\_With\_DB/})} \\
\midrule
\endhead
Embedding model &
\texttt{all-MiniLM-L6-v2}, 384-dim, \texttt{normalize\_embeddings=True}
(line 375) &
\texttt{all-MiniLM-L6-v2}, 384-dim, normalised by the model's own
\texttt{Normalize} module (backend line 22) \\
\addlinespace
Generator LLM &
\texttt{claude-sonnet-4-6} via the \texttt{anthropic} SDK (line 562) &
\texttt{openai/gpt-oss-120b} via \textbf{Groq} (backend line 21) \\
\addlinespace
Chunking &
Character window, 800 size / 150 overlap, snapped backwards to the last space in
the second half of the window (lines 181--317) &
\texttt{RecursiveCharacterTextSplitter}, 1000 / 200; separate no-overlap cue
grouper for \texttt{.srt}/\texttt{.vtt} (loader lines 8, 29--55) \\
\addlinespace
Index &
One \texttt{(N, 384)} NumPy array plus an aligned list of dicts (lines 404--414) &
Chroma (HNSW), persisted per transcript to
\texttt{storage/chroma\_store} (backend lines 184--200) \\
\addlinespace
Search &
\texttt{matrix @ q} then \texttt{argsort}, exact, $k=4$ (lines 450--464) &
\texttt{as\_retriever(search\_kwargs=\{"k": 4\})}, approximate
(backend line 231) \\
\addlinespace
Input formats &
\texttt{.txt} and \texttt{.md} only (line 79) &
\texttt{.txt}, \texttt{.pdf}, \texttt{.srt}, \texttt{.vtt} (frontend lines 96--99) \\
\addlinespace
Citations &
\textbf{Yes.} Prompt asks for \texttt{[2]}-style chunk citations and the CLI
prints retrieval scores before every answer (lines 517--519, 603--604) &
\textbf{No.} Timestamps go into the prompt, but the UI renders answer text only
(frontend lines 183--193) \\
\addlinespace
Conversation memory &
None. Single-turn; each question is independent (lines 588--606) &
\texttt{SqliteSaver} checkpointer, resumable named threads
(backend lines 31--32, 254) \\
\addlinespace
Persistence &
None. Re-embeds the whole corpus at every start (line 414) &
Three tiers: Chroma on disk, checkpoint DB, \texttt{transcript\_meta} table \\
\addlinespace
Interface &
Terminal REPL (lines 588--606) &
Streamlit chat UI with token streaming \\
\bottomrule
\end{longtable}

\noindent
\textbf{Worth noticing, because it inverts the usual story:} the from-scratch
version is the one that cites its sources and shows retrieval scores. The
framework version has more capability --- persistence, multi-turn memory, PDF and
subtitle parsing, streaming --- but strictly \emph{worse} answer provenance. That
is a genuinely good thing to say out loud in an interview, and it is the kind of
observation that only comes from having built both.

% =====================================================================
\section{Head-to-head benchmark}\label{sec:bench}
% =====================================================================

\subsection{Method}

\begin{itemize}
  \item \textbf{Corpus.} The AI-infrastructure transcript this application
        genuinely indexed, recovered from the live Chroma collection and
        de-overlapped back into 52{,}826 characters of raw text
        (\texttt{eval/ai\_infra\_transcript.txt}).
  \item \textbf{Questions.} 35, spanning the whole transcript
        (\texttt{eval/eval\_questions.json}).
  \item \textbf{Gold labels.} Each question carries an \texttt{answer\_span}; a
        chunk is gold if it contains that span. Labels are therefore recomputed
        for each chunking scheme, which is the only fair way to compare systems
        that chunk differently. It also handles overlap correctly: a fact
        straddling a boundary is legitimately present in two chunks.
  \item \textbf{Fairness.} Both pipelines receive the \emph{same raw document}
        and chunk it themselves, because chunking is part of what is being
        compared. The NumPy pipeline is driven through its own
        \texttt{split()}, \texttt{embed()} and \texttt{retrieve()} functions ---
        imported, not reimplemented. The encoder is warm in both, and each query
        is timed three times with the minimum taken.
\end{itemize}

\subsection{Result at shipped defaults}

\begin{center}\small
\begin{tabular}{@{}lrr@{}}
\toprule
\textbf{Metric} & \textbf{Framework (1000/200)} & \textbf{From-scratch (800/150)} \\
\midrule
Chunks produced       & 65 (943 chars avg) & 82 (792 chars avg) \\
hit@1                 & 0.486 & 0.486 \\
hit@3                 & 0.743 & \textbf{0.771} \\
\textbf{hit@4}        & 0.829 & \textbf{0.857} \\
hit@5                 & 0.829 & \textbf{0.886} \\
hit@10                & 0.943 & \textbf{0.971} \\
recall@4              & 0.781 & \textbf{0.814} \\
MRR                   & 0.633 & \textbf{0.643} \\
\midrule
Retrieval latency, mean & 9.6\,ms & \textbf{7.9\,ms} \\
Retrieval latency, p95  & 10.6\,ms & \textbf{8.6\,ms} \\
Index build (warm encoder) & 1.25\,s & \textbf{0.94\,s} \\
\bottomrule
\end{tabular}
\end{center}

\noindent
Read literally: the 190-line NumPy script retrieves slightly better and slightly
faster than the framework pipeline. Stopping there would be a mistake, and an
interviewer will push on it --- so the next subsection is the important one.

\subsection{Ablation: is it the implementation, or the chunk size?}

The two systems differ in two ways at once, which is a confound. Holding chunk
size fixed and re-running both isolates it. \textbf{hit@4} across the grid:

\begin{center}\small
\begin{tabular}{@{}lrrl@{}}
\toprule
\textbf{Chunk size / overlap} & \textbf{Framework} & \textbf{From-scratch} & \textbf{Difference} \\
\midrule
500 / 100   & 0.800 & 0.714 & framework +0.086 (3 questions) \\
800 / 150   & \textbf{0.886} & 0.857 & framework +0.029 (1 question) \\
1000 / 200  & 0.829 & 0.857 & from-scratch +0.028 (1 question) \\
\bottomrule
\end{tabular}
\end{center}

\noindent
And the same grid on MRR: framework 0.642 / \textbf{0.700} / 0.633, from-scratch
0.567 / 0.643 / 0.639.

\paragraph{What this establishes.}
\begin{enumerate}
  \item \textbf{Chunk size dominates.} Within the from-scratch pipeline alone,
        moving from 500 to 800 characters moves hit@4 by 14 points (5 questions).
        The largest implementation gap at a fixed size is 3 points (1 question).
  \item \textbf{The default the application ships is not its best setting.}
        The framework pipeline at 800/150 scores hit@4 0.886 and hit@10
        \textbf{1.000} --- every question answerable within ten chunks --- against
        0.829 and 0.943 at the 1000/200 it currently uses. That is a concrete,
        free improvement to the shipped app.
  \item \textbf{The original ``NumPy wins'' reading was an artefact.} It compared
        800/150 against 1000/200 and attributed a chunking effect to the
        implementation.
\end{enumerate}

\paragraph{Honesty about the sample size --- say this before you are asked.}
There are 35 questions, so one question is worth 2.9 points of hit@4. Differences
of one or two questions (every implementation gap in the table above) are
\textbf{within noise} and should not be claimed as real. Only the chunk-size
effects, which move 3--5 questions, are large enough to lean on. A second honest
caveat: because gold labels are recomputed per chunking, the average number of
gold chunks per question varies from 1.11 to 1.31 across configurations, which
slightly favours the configurations with more gold chunks. Both caveats are worth
stating in an interview; volunteering them reads as rigour, while being caught
without them reads as overclaiming.

\subsection{What each side actually buys}

\begin{longtable}{@{}R{4.6cm}R{10.6cm}@{}}
\toprule
\textbf{The framework earns} & \textbf{The evidence} \\
\midrule
\endhead
A better splitter &
Recursive splitting on paragraph and sentence boundaries beats naive
space-snapping at 500 and 800 characters. \\
Capability, not quality &
PDF and subtitle parsing, persistence, checkpointed multi-turn threads, token
streaming --- none of which the NumPy version has, and none of which this
benchmark measures. \\
\midrule
\textbf{The framework costs} & \\
\midrule
Latency & 9.6\,ms versus 7.9\,ms per query; roughly 20\% overhead for a wrapper
around the same encoder and the same arithmetic. \\
Nothing on quality at this scale &
65 vectors is far below where an approximate index matters; exact NumPy search
returns the same neighbours. \\
Visibility &
The retrieval scores that \texttt{rag\_annotated.py} prints on every query --- the
single best debugging signal in a RAG system --- are not surfaced anywhere in
the framework app. \\
\bottomrule
\end{longtable}

% =====================================================================
\section{Verdicts on the original three bullets}
% =====================================================================

\subsection{Bullet 1 --- from-scratch NumPy}\label{sec:bullet1}

\begin{quote}\itshape
Built a RAG pipeline from scratch in NumPy, implementing document chunking,
embedding, vector similarity search and context assembly directly, to establish
the underlying mechanics without framework abstraction.
\end{quote}

\begin{longtable}{@{}R{6.0cm}R{2.4cm}R{6.6cm}@{}}
\toprule
\textbf{Claim} & \textbf{Verdict} & \textbf{Evidence} \\
\midrule
\endhead
Chunking implemented directly & \SUP &
Hand-written sliding window with overlap and word-boundary snapping, including a
guard against the degenerate \texttt{overlap >= size} case (lines 181--317). \\
\addlinespace
Vector similarity search implemented directly & \SUP &
\texttt{matrix @ q} plus \texttt{argsort}, with normalisation at embed time so
cosine reduces to a dot product (lines 375, 450--464). \\
\addlinespace
Context assembly implemented directly & \SUP &
Numbered, source-tagged context blocks joined for the prompt (lines 546--548). \\
\addlinespace
``Embedding'' implemented directly & \PART &
\textbf{This is the one word to fix.} The encoder is pretrained
\texttt{all-MiniLM-L6-v2} loaded from sentence-transformers (line 354). You wrote
the pipeline around it; you did not implement an embedding model. \\
\addlinespace
``From scratch'' & \PART &
True of the retrieval machinery, not of the encoder. Your own phrasing --- ``not
totally from scratch, but the chunking and retrieval are mine in NumPy'' --- is
the accurate one and should go on the resume in that spirit. \\
\bottomrule
\end{longtable}

\subsection{Bullet 2 --- the framework assistant}

Unchanged from version 1: \SUP{} on LangChain, LangGraph, Streamlit, local
serving, upload, and semantic retrieval; \PART{} on ``grounded'' (prompt-level
only, no threshold, no citation rendering, no faithfulness test) and on
``document assistant'' (it is specifically a \emph{transcript} assistant).

\subsection{Bullet 3 --- the comparison}

Now \SUP, for the first time: Section~\ref{sec:bench} supplies retrieval quality,
latency, indexing cost, an ablation isolating the cause, and a statement of where
the abstraction helps and where it hides things. Bullet 3 was the weakest line on
the resume and is now the strongest.

\subsection{The two words that still fail}

\textbf{``Deployed''} --- no Dockerfile, no host config, no public URL, and a
\texttt{requirement.txt} that omits Streamlit, LangGraph, Chroma, Groq and
sentence-transformers. The app runs under \texttt{streamlit run} on localhost.
\textbf{``Agent''} --- two nodes, three unconditional edges, no tool binding, no
conditional routing; the LLM never influences control flow.

% =====================================================================
\section{Resume: title and bullets}
% =====================================================================

\subsection{Project title}

Drop ``Deployed Document Assistant''. In order of preference:

\begin{enumerate}
  \item \textbf{Retrieval-Augmented Generation: From-Scratch and Framework
        Implementations, Benchmarked} \\
        Leads with the comparison, which is what makes the project unusual.
  \item \textbf{Retrieval-Augmented Generation: NumPy Pipeline and LangGraph
        Transcript Assistant} \\
        Plainer; leads with the artefacts.
  \item \textbf{RAG from First Principles: Building, Wrapping, and Measuring a
        Retrieval Pipeline} \\
        Strongest if the reader is technical; slightly informal for a screening
        filter.
\end{enumerate}

\noindent
Dates: Aug 2026 -- Present.

\subsection{Primary bullets (use these three)}

\begin{enumerate}[leftmargin=*,itemsep=5pt]

\item Implemented a RAG retrieval pipeline in NumPy around a MiniLM-384d encoder
--- overlapping character chunking, unit-normalised embeddings, dot-product cosine
search over an $(N \times 384)$ matrix, and cited context assembly --- to work
without retrieval-framework abstractions.

\item Built a transcript QA assistant on LangChain/LangGraph/Streamlit (Chroma,
Groq \texttt{gpt-oss-120b}): a two-node retrieve--generate graph with
SQLite-checkpointed resumable chat threads, \texttt{txt}/\texttt{pdf}/\texttt{srt}/\texttt{vtt}
ingestion, and map-reduce summarisation under a token-rate cap.

\item Benchmarked both on a 35-question span-labelled gold set over a
53k-character corpus: hit@4 0.86 vs 0.83 at 8--10\,ms retrieval, and a chunk-size
ablation isolating chunking --- not the framework --- as the dominant factor,
lifting hit@4 to 0.89 and hit@10 to 1.00.

\end{enumerate}

\noindent
\textbf{On length:} at normal resume font each of these runs to two lines. If you
need three single lines, use the compressed set below and keep the detail for the
interview.

\subsection{Compressed variants (one line each)}

\begin{enumerate}[leftmargin=*,itemsep=4pt]
  \item Implemented RAG retrieval in NumPy --- overlap chunking, normalised
        MiniLM-384d embeddings, dot-product cosine search --- with no retrieval framework.
  \item Built a LangGraph/Streamlit transcript QA assistant: two-node
        retrieve--generate graph, Chroma store, SQLite-checkpointed resumable threads, 4 input formats.
  \item Benchmarked both on 35 span-labelled questions: hit@4 0.86 vs 0.83;
        ablation showed chunk size, not the framework, drives retrieval quality (hit@10 1.00 at 800/150).
\end{enumerate}

\subsection{Emphasis by target role}

\begin{longtable}{@{}R{4.0cm}R{11.2cm}@{}}
\toprule
\textbf{Role} & \textbf{What to foreground} \\
\midrule
\endhead
Capital One --- Associate, Business Analysis &
Lead with bullet 3. The story is measurement and a controlled experiment that
overturned a first reading: you designed a gold set, found a confound, ran an
ablation, and reported effect sizes against sample-size caveats. Keep the
implementation bullets short --- they are context, not the point. \\
\addlinespace
Mastercard --- Data Scientist II &
Lead with bullet 3, then bullet 1. Emphasise evaluation design (span-derived
labels, hit@k/MRR, per-configuration relabelling) and the explicit statement that
one-question differences at $n=35$ are noise. That sentence is the most senior
thing on the page. \\
\addlinespace
EXL --- AI Engineer &
Lead with bullet 2, then 3. Foreground the system: LangGraph state and
checkpointing, persistence tiers, rate-limit handling with map-reduce and
backoff, streaming. Then use bullet 3 to show you measure what you build. \\
\bottomrule
\end{longtable}

\noindent
\textbf{Vocabulary to keep consistent everywhere:} \emph{served locally} (not
deployed), \emph{two-node graph} / \emph{pipeline} (not agent),
\emph{context-constrained} (not verified-grounded), \emph{retrieval pipeline in
NumPy} (not from-scratch embeddings).

% =====================================================================
\section{Interview questions this version invites}
% =====================================================================

Version 1 listed ten questions about the framework application; they all still
apply. These five are new, and specific to the comparison.

\begin{enumerate}[leftmargin=*,itemsep=6pt]

\item \textbf{Your NumPy version beat the framework. Why not just delete the
framework one?}\\
\emph{Answer:} Because the win was not real. At the defaults it was 0.857 versus
0.829, but the two systems chunked differently --- 800/150 against 1000/200. When I
held chunk size constant the gap reversed or vanished, and at 800/150 the
framework pipeline was actually best at hit@4 0.886 and hit@10 1.000. The
difference was chunk size, not NumPy versus Chroma. What the framework really
buys is capability --- PDF and subtitle parsing, persistence, checkpointed
multi-turn threads --- none of which my benchmark measures.

\item \textbf{Your benchmark says hit@4 0.886 versus 0.857. Is that a real
difference?}\\
\emph{Answer:} No. With 35 questions, one question is 2.9 points, so that gap is a
single question and well inside noise. The only effects I would defend are the
chunk-size ones, which move three to five questions. If I needed the
implementation comparison to be conclusive I would need several hundred questions
across multiple documents, or a paired test on per-question outcomes.

\item \textbf{How did you build the gold labels without hand-annotating
everything?}\\
\emph{Answer:} Each question carries a short answer span, and any chunk containing
that span is gold. It makes labelling reproducible and auditable rather than a
list of indices I typed, and it survives re-chunking --- which mattered, because
the ablation changes chunk boundaries and I needed labels recomputed for each
configuration. It also handles overlap honestly: a fact straddling a boundary is
genuinely in two chunks and both count.

\item \textbf{You compared retrieval, but RAG is retrieval plus generation. What
did you not measure?}\\
\emph{Answer:} Generation, and I am explicit about it. I did not measure answer
faithfulness, refusal behaviour on unanswerable questions, or end-to-end latency,
because those need live API calls against two different providers --- Anthropic in
one pipeline, Groq in the other --- and are confounded by the model choice rather
than the retrieval design. Retrieval is also the half worth fixing first: if the
right chunk is not retrieved, no model recovers.

\item \textbf{Which implementation would you take to production?}\\
\emph{Answer:} The framework one, at 800/150 rather than its current 1000/200, and
I would port two things from the NumPy version: the retrieval scores, which are
the best debugging signal in the system, and the numbered chunk citations, so an
answer can be checked against its source. Ironically the 190-line script has
better provenance than the application does.

\end{enumerate}

% =====================================================================
\section{What to do next}
% =====================================================================

\begin{enumerate}[leftmargin=*]
  \item \textbf{Rotate the five API keys still committed in \texttt{.env}.}
        Unchanged from version 1, and still the most urgent item. The file is in
        the initial commit and the repository has a GitHub remote.
  \item \textbf{Change the application's splitter to 800/150.} One line, measured,
        and worth 5.7 points of hit@4 plus a perfect hit@10 on this corpus.
  \item \textbf{Surface retrieval scores and citations in the Streamlit UI.} Port
        what \texttt{rag\_annotated.py} already does. It makes ``grounded''
        defensible rather than aspirational.
  \item \textbf{Stop tracking the \texttt{langchain/} virtual environment in git}
        (18{,}110 of 18{,}184 changed files) and write a
        \texttt{requirements.txt} that can actually install the app.
  \item \textbf{Optional, if a live URL matters:} deploy on Hugging Face Spaces,
        which has the RAM for \texttt{torch} and purchasable persistent storage
        for the Chroma and checkpoint databases. Only then does the word
        ``deployed'' belong on the resume.
\end{enumerate}

\vspace{0.8em}
\noindent
\textbf{Closing note.} The strongest thing here is no longer either
implementation. It is that you built the same system twice, measured both against
one gold set, found a confound in your own first result, ran an ablation to
isolate it, and can state the effect sizes alongside what the sample size does
and does not support. Most candidates present a project that works. Presenting a
project you have genuinely interrogated is a different conversation.

\end{document}
